\begin{table}[htbp]
\centering
\footnotesize
\begin{threeparttable}
\caption{Representative recovery mechanisms and where each one pays its main cost.}
\label{tab:mechanism-summary}
\begin{tabularx}{\linewidth}{@{}p{0.23\linewidth}p{0.15\linewidth}X p{0.21\linewidth}@{}}
\toprule
\textbf{Mechanism family} & \textbf{Main resource} & \textbf{What it buys} & \textbf{Transfer boundary} \\
\midrule
Dense MAP baseline & R1, R7, R8 & Cheap baseline behavior and exact recovery on easy cells. & Capacity-limited on harder cells. \\
Native BCF substrate & R4, R12 & Strong clean common-envelope recovery under a substrate designed for the task. & Different native contract from dense MAP. \\
Linear-code recovery & R4, R5 & Exact recovery when explicit GF(2) code constraints are part of the task. & Narrow algebraic envelope. \\
Context-conditioned search & R6, sometimes R13 & Candidate-space narrowing before expensive recovery. & Better local recall does not guarantee a better end-to-end operating point. \\
Exact sidecar replay & R5, R14 & Safe deterministic replay after retrieval when exact structure is already known. & Does not solve initial localization by itself. \\
Soft-information repair & R2, R3, R11 & Preservation of local reliability signal or confidence information. & Often loses to simpler equal-bit controls after full accounting. \\
\bottomrule
\end{tabularx}
\begin{tablenotes}[flushleft]
\footnotesize
\item Reader-facing labels are used in the main paper; the supplementary atlas provides the full mapping to normalized repository evidence records.
\end{tablenotes}
\end{threeparttable}
\end{table}
